% Подключаем все нужные пакеты
%%% Локализация и шрифты
\usepackage{asymptote}
\usepackage[utf8]{inputenc}          % Кодировка UTF-8
\usepackage[T2A]{fontenc}            % Русская кодировка
\usepackage[english,russian]{babel} % Локализация
\usepackage{cmap}                    % Поддержка поиска в PDF
\usepackage{tikz}
\usepackage{xstring}
\usepackage[left=20mm, top=10mm, right=20mm, bottom=10mm, nohead, nofoot]{geometry}
\renewcommand{\baselinestretch}{1.0}% Межстрочный интервал
\usepackage{ragged2e}               % Выравнивание текста
\usepackage{microtype}              % Улучшение качества вывода текста
\usepackage{multicol}
\usepackage{tcolorbox}
\tcbuselibrary{listingsutf8, breakable, skins}
\justifying
\sloppy
\tolerance=500
\hyphenpenalty=10000
\emergencystretch=3em

\usepackage[normalem]{ulem}
\pagestyle{empty} \textwidth=19.0cm \oddsidemargin=-1.3cm
\textheight=26cm \topmargin=-3.0cm
\tcbuselibrary{listingsutf8}
\usepackage{pgfplots}
\renewcommand{\rmdefault}{cmss}
\usepgfplotslibrary{patchplots}
\pgfplotsset{compat=1.18}
\usetikzlibrary{fillbetween}
\usetikzlibrary{patterns}
% Создаём новый счётчик задач

% определение
\newcounter{maindf}
\newcounter{subdf}[maindf]

\newcommand{\dfsub}{\arabic{maindf}.\arabic{subdf}}

% теорема
\newcounter{mainth}
\newcounter{subth}[mainth]

\newcommand{\thsub}{\arabic{mainth}.\arabic{subth}}

% утверждение
\newcounter{mainpr}
\newcounter{subpr}[mainpr]

\newcommand{\propsub}{\arabic{mainpr}.\arabic{subpr}}

% пример
\newcounter{mainex}
\newcounter{subex}[mainex]

\newcommand{\exsub}{\arabic{mainex}.\arabic{subex}}

% задача
\newcounter{mainta}
\newcounter{subta}[mainta]

\newcommand{\tasub}{\arabic{mainta}.\arabic{subta}}

\newcounter{mainle}
\newcounter{suble}[mainle]

\newcommand{\lesub}{\arabic{mainle}.\arabic{suble}}

\renewcommand{\familydefault}{\ttdefault}

% Определяем стили блоков
\usepackage{xcolor}
% myrose = #FFB5C1
% myred = #FF6464
% mybrown = #D2B48C
% mycyan = #87CEEB
% myorange = #FFB000
% mypurple = #D8BFD8
% mygreen = #14DE14
% --- Однострочное создание блоков \problem, \solution, \answer ---
\definecolor{myyellow}{HTML}{737700} % Жёлтый
\definecolor{mygreen}{HTML}{03810b}  % Голубой
\definecolor{mycyan}{HTML}{87CEEB}  % Голубой

\newtcolorbox{SectionBox}{
  breakable,
  enhanced,
  colback=yellow,
  colframe=myyellow,
  fonttitle=\ttfamily,
}

\newcommand{\sect}[1]{%
  \refstepcounter{subdf}%
  \begin{SectionBox}%
    \begin{center}
      #1
    \end{center}%
  \end{SectionBox}%
}

\newtcolorbox{SubsectionBox}{
  breakable,
  enhanced,
  colback=green,
  colframe=mygreen,
  fonttitle=\ttfamily,
}

\newcommand{\subsect}[1]{%
  \refstepcounter{subdf}%
  \begin{SubsectionBox}%
    \begin{center}
      #1
    \end{center}%
  \end{SubsectionBox}%
}

\newtcolorbox{SubsubsectionBox}{
  breakable,
  enhanced,
  colback=mycyan,
  colframe=cyan,
  fonttitle=\ttfamily,
}

\newcommand{\subsubsect}[1]{%
  \refstepcounter{subdf}%
  \begin{SubsubsectionBox}%
    \begin{center}
      #1
    \end{center}%
  \end{SubsubsectionBox}%
}

% Команда для векторов
\newcommand{\vect}[1]{\mathbf{#1}}
\usepackage[upint]{stix2}

%%% Колонтитулы
\usepackage{fancyhdr}
\pagestyle{fancy}
\renewcommand{\headrulewidth}{0.3mm} % Толщина линии в колонтитулах

\rhead{Динамическая оптимизация в экономике и финансах}
\lhead{ФЭН ВШЭ, 2026}
\headsep=25pt
\footskip=15mm
\textheight=700pt
\headheight=35pt % Возможно, понадобится увеличить до ~30pt. Проверьте log-файл.
\parindent=0mm
%%% Математика
\usepackage{amsmath, amsfonts, amssymb, amsthm, mathtools} % Базовые математические пакеты
\usepackage{icomma}   
\usepackage{euscript}               % Шрифт Евклид
\usepackage{mathrsfs}               % Красивый матшрифт
\usepackage{cancel}                 % Зачёркивание в формулах
\usepackage{diagbox}
\usepackage[unicode, colorlinks, urlcolor=darkblue, linkcolor=violet]{hyperref}
\usetikzlibrary{matrix}
% Замены часто употребляемых шрифтов
\newcommand{\mb}[1]{\mathbb{#1}}
\newcommand{\mc}[1]{\mathcal{#1}}
\newcommand{\ms}[1]{\mathscr{#1}}

\newcommand{\cov}{\operatorname{cov}}
\newcommand{\Var}{\operatorname{Var}}
\newcommand{\corr}{\operatorname{corr}}
\newcommand{\const}{\operatorname{const}}
\newcommand{\Be}{\operatorname{Be}}
\newcommand{\Bi}{\operatorname{Bi}}
\newcommand{\Poi}{\operatorname{Poi}}
\newcommand{\Exp}{\operatorname{Exp}}
\newcommand{\Geo}{\operatorname{Geo}}
\newcommand{\bias}{\operatorname{bias}}
\newcommand{\MSE}{\operatorname{MSE}}
\newcommand{\sign}{\operatorname{sign}}
\newcommand{\plim}{\operatorname{plim}}
\newcommand{\asi}{\overset{asy}{\sim}}
\newcommand{\argmax}[1]{\underset{#1}{\operatorname{argmax}}}